\clearpage
\section{Follow one real execution from the command to the recovered key}
This section is an execution route, not another independent description of the
functions. Read the numbered stages in order: the output of one call becomes
an input to the next. Numerical records are generated from the current solver.
The builder also runs the real terminal entry point without interactive pauses
and verifies that it succeeds. The command below adds pauses but follows the
same mathematical path:
\begin{Verbatim}
python src/piret_sm4.py --faults 8 --row 0 --step
\end{Verbatim}

\subsection{The route through the files}
\begin{center}
\begin{tikzpicture}[node distance=5mm,box/.style={draw,rounded corners,fill=teal!5,text width=125mm,align=center,minimum height=11mm,font=\small},flow/.style={-{Stealth},thick}]
\node[box] (a) {\hx{piret\_sm4.py}: entry guard $\to$ logger $\to$ \hx{main()}};
\node[box,below=of a] (b) {\hx{run.py}: \hx{sm4\_demo(terminal, limit=8, start\_byte=0)}};
\node[box,below=of b] (c) {Target: \hx{sm4\_fault\_encrypt(P, K, byte=None)} $\to$ clean ciphertext};
\node[box,below=of c] (d) {Solver: \hx{recover\_sm4\_master\_key(C, oracle, masks=1..8)}};
\node[box,below=of d] (e) {Per observation: peel $\to$ oracle/target $\to$ peel $\to$ byte filter $\to$ intersection};
\node[box,below=of e] (f) {After four subkeys: \hx{reverse\_sm4()} $\to$ master key $\to$ verification};
\draw[flow] (a)--(b);\draw[flow] (b)--(c);\draw[flow] (c)--(d);
\draw[flow] (d)--(e);\draw[flow] (e)--(f);
\end{tikzpicture}
\end{center}
The target knows the secret key. The solver receives ciphertexts and fault
descriptors. The terminal's teaching callbacks display additional internals;
their displays are not secret inputs to the candidate equations.

\subsection{Stage 1: argument parsing and the clean target call}
The entry guard calls \hx{run\_with\_log(main, "piret\_sm4")}. The parser obtains
\hx{faults=8}, \hx{row=0}, \hx{step=True}, with automatic color mode.
\hx{main()} constructs \hx{LessonTerminal} and invokes \hx{run.sm4\_demo}.
That function converts \hx{sm4.KEY} and \hx{sm4.PLAINTEXT} from hexadecimal
strings to bytes. For this public demonstration vector:
\[
P=K=\hx{0123456789ABCDEFFEDCBA9876543210}.
\]
The clean call uses \hx{byte=None}; the observer records events without
injecting any fault. The result is
\[
C=\hx{681EDF34D206965E86B3E94F536E4246}.
\]
The teacher can display the clean round-31 trace
$A=\hx{96FF9D45}$, $B=\hx{1A488C73}$, $L(B)=\hx{138D5078}$.
These values explain the target; the solver does not receive this trace.

\subsection{Stage 2: enter the solver and select the first byte}
The actual call passes \hx{range(1, 9)} as permitted masks and the two
callbacks \hx{oracle} and \hx{report}. Solver state starts as
\begin{Verbatim}
keys = {}
history = []
counts = {}
used = 0
\end{Verbatim}
The outer round loop selects \hx{number=31}. The call
\hx{peel\_sm4(C, 31, \{\})} undoes no rounds and returns $C$ unchanged.
The subkey buffer is \hx{00 00 00 00}; \hx{start\_byte=0} produces visiting
order $[0,1,2,3]$. The first selected byte is 0, \hx{survivors=None}, and
\hx{enumerate(masks, 1)} produces observation 1, mask 01.

\subsection{Stage 3: the oracle enters the target, then returns only ciphertext}
\begin{center}
\begin{tikzpicture}[x=1mm,y=1mm,every node/.style={font=\small}]
\node[draw,fill=teal!8,minimum width=38mm,minimum height=11mm] (s) at (0,0) {solver};
\node[draw,fill=blue!8,minimum width=58mm,minimum height=11mm] (t) at (98,0) {oracle $\to$ target};
\draw[-{Stealth},thick] (s.north east) to[bend left=18] node[above] {round 31, byte 0, mask 01} (t.north west);
\draw[-{Stealth},thick] (t.south west) to[bend left=18] node[below] {faulty ciphertext only} (s.south east);
\end{tikzpicture}
\end{center}
\hx{oracle(31, 0, 1)} calls
\hx{sm4\_fault\_encrypt(P, K, byte=0, mask=1, fault\_round=31)}.
The callback intervenes only at the selected S-box input:
\[
\begin{aligned}
\text{shift}&=24-8(0)=24,\\
\text{shifted mask}&=\hx{01}\ll24=\hx{01000000},\\
A'&=\hx{96FF9D45}\xor\hx{01000000}=\hx{97FF9D45}.
\end{aligned}
\]
After the target completes normal SM4 processing, the oracle returns
\[
F_1=\hx{1B50E209D206965E86B3E94F536E4246}.
\]
The terminal may illustrate fault propagation using its teacher trace.
The solver gets $F_1$ alone. It calls \hx{peel\_sm4(F1,31,keys)}, which also
undoes no rounds.

\subsection{Stage 4: follow the first byte-filter call}
\hx{sm4\_byte\_candidates(C,F1,byte=0,mask=1)} unpacks four big-endian words
and checks that the last three words match. It computes
\[
\Delta X=\hx{681EDF34}\xor\hx{1B50E209}=\hx{734E3D3D}.
\]
Its call to \hx{inverse\_linear(0x734E3D3D)} obtains the fixed inverse rows.
On a fresh process, \hx{inverse\_rows()} builds and caches them; subsequent
observations reuse the matrix. The earlier illustrated matrix sections explain
that preparation. Applying its rows gives
\[
\Delta B=\hx{4F000000},\qquad
d=(\hx{4F000000}\gg24)\mathbin{\&}\hx{FF}=\hx{4F}.
\]
The other three words yield the public constant
\[
q_{\rm word}=\hx{D206965E}\xor\hx{86B3E94F}\xor\hx{536E4246}
=\hx{07DB3D57},\qquad q=\hx{07}.
\]
The 256-input loop retains only $a=\hx{96}$ and $a=\hx{97}$ because
$S(a)\xor S(a\xor\hx{01})=\hx{4F}$.
Converting input guesses to key guesses gives
$\hx{96}\xor\hx{07}=\hx{91}$ and
$\hx{97}\xor\hx{07}=\hx{90}$.
The function returns
\begin{Verbatim}
fresh  = {0x90, 0x91}
inputs = {0x96, 0x97}
delta  = 0x4F000000
\end{Verbatim}
Back in the solver, \hx{survivors=None} selects the first-observation branch:
\hx{before=256} and \hx{survivors=set(fresh)}.
\hx{used} becomes 1 and \hx{counts["31:0"]} becomes 2. An event is appended
and passed to \hx{report(event)} for display. Since two candidates remain,
the solver does not write a subkey byte yet.

\subsection{Stage 5: the second observation resolves the first byte}
The next observation requests \hx{oracle(31,0,2)}. It returns
\[
F_2=\hx{ED9DD932D206965E86B3E94F536E4246}.
\]
The same filter computes
\[
\begin{aligned}
\Delta X&=\hx{681EDF34}\xor\hx{ED9DD932}=\hx{85830606},\\
\Delta B&=L^{-1}(\hx{85830606})=\hx{81000000},\qquad d=\hx{81}.
\end{aligned}
\]
Its compatible S-box inputs are $\{\hx{75},\hx{77},\hx{94},\hx{96}\}$,
giving fresh key guesses $\{\hx{70},\hx{72},\hx{93},\hx{91}\}$.
The existing survivor set is now intersected:
\[
\{\hx{90},\hx{91}\}\cap\{\hx{70},\hx{72},\hx{93},\hx{91}\}
=\{\hx{91}\}.
\]
Now \hx{used=2}, \hx{counts["31:0"]=1} and \hx{subkey[0]=0x91}.
The buffer becomes \hx{91 00 00 00}; the other positions are not recovered yet.
The \hx{break} exits the mask loop, so mask 03 is not requested for this byte.

\subsection{Stages 6--9: follow all observations and assemble each subkey}
The next tables are generated from the real event history, in execution order.
They connect loop variables, inverse differences, candidate intersections,
byte-buffer updates and subkey milestones. All hexadecimal entries denote
values, not variable names. Query numbering excludes the clean reference.
\subsubsection*{Round 31: input views and the byte loop}
Before this round, recovered subkeys are:
\hx{none}.
The clean view returned by \texttt{peel\_sm4} is
\[C_{\rm view}=\hx{681EDF34D206965E86B3E94F536E4246}.\]
The first oracle response for this round and its peeled view are
\[\begin{aligned}F&=\hx{1B50E209D206965E86B3E94F536E4246},\\F_{\rm view}&=\hx{1B50E209D206965E86B3E94F536E4246}.\end{aligned}\]
The filter operates on these views. In the table, ``fresh'' is the key-byte set for this pair and ``remaining'' is the intersection with earlier observations for the same byte.
\begin{center}\small\begin{tabular}{rrrlll}\toprule
Query & Byte & Mask & $\Delta B$ & Fresh & Remaining\\\midrule
1 & 0 & 01 & 4F000000 & 90 91 & 90 91\\
2 & 0 & 02 & 81000000 & 70 72 91 93 & 91\\
3 & 1 & 01 & 00710000 & 24 25 & 24 25\\
4 & 1 & 02 & 00830000 & 24 26 & 24\\
5 & 2 & 01 & 00007900 & A0 A1 & A0 A1\\
6 & 2 & 02 & 00006F00 & A0 A2 & A0\\
7 & 3 & 01 & 00000080 & 12 13 & 12 13\\
8 & 3 & 02 & 000000C9 & 10 12 & 12\\
\bottomrule\end{tabular}\end{center}
Unique bytes are written into the subkey buffer in the following order. Zero-filled positions not yet written are placeholders, not recovered bytes.
\begin{center}\small\begin{tabular}{rrlr}\toprule
Byte written & Value & Buffer after write & Queries used\\\midrule
0 & 91 & 91000000 & 2\\
1 & 24 & 91240000 & 4\\
2 & A0 & 9124A000 & 6\\
3 & 12 & 9124A012 & 8\\
\bottomrule\end{tabular}\end{center}
Only now is \texttt{keys[31]} assigned: $rk_{31}=\hx{9124A012}$. This word may be used to peel the next attacked round.
\subsubsection*{Round 30: input views and the byte loop}
Before this round, recovered subkeys are:
\hx{rk31=9124A012}.
The clean view returned by \texttt{peel\_sm4} is
\[C_{\rm view}=\hx{D206965E86B3E94F536E42467B938F4C}.\]
For the clean ciphertext, the peeling calls perform the following actual calculations in order. The last column is the recovered older state word:
\begin{center}\small\begin{tabular}{rllll}\toprule
Undone round & S-box input & $\tau$ output & $L$ output & Older word\\\midrule
31 & 96FF9D45 & 1A488C73 & 138D5078 & 7B938F4C\\
\bottomrule\end{tabular}\end{center}
The first oracle response for this round and its peeled view are
\[\begin{aligned}F&=\hx{AE21A70A772712DA86B3E94F536E4246},\\F_{\rm view}&=\hx{772712DA86B3E94F536E42467B938F4C}.\end{aligned}\]
The filter operates on these views. In the table, ``fresh'' is the key-byte set for this pair and ``remaining'' is the intersection with earlier observations for the same byte.
\begin{center}\small\begin{tabular}{rrrlll}\toprule
Query & Byte & Mask & $\Delta B$ & Fresh & Remaining\\\midrule
9 & 0 & 01 & 21000000 & 00 01 & 00 01\\
10 & 0 & 02 & 3C000000 & 01 03 & 01\\
11 & 1 & 01 & 00550000 & CE CF & CE CF\\
12 & 1 & 02 & 006D0000 & CD CF & CF\\
13 & 2 & 01 & 00005100 & 72 73 & 72 73\\
14 & 2 & 02 & 0000AB00 & 70 72 & 72\\
15 & 3 & 01 & 000000EB & E4 E5 & E4 E5\\
16 & 3 & 02 & 000000FF & E5 E7 & E5\\
\bottomrule\end{tabular}\end{center}
Unique bytes are written into the subkey buffer in the following order. Zero-filled positions not yet written are placeholders, not recovered bytes.
\begin{center}\small\begin{tabular}{rrlr}\toprule
Byte written & Value & Buffer after write & Queries used\\\midrule
0 & 01 & 01000000 & 10\\
1 & CF & 01CF0000 & 12\\
2 & 72 & 01CF7200 & 14\\
3 & E5 & 01CF72E5 & 16\\
\bottomrule\end{tabular}\end{center}
Only now is \texttt{keys[30]} assigned: $rk_{30}=\hx{01CF72E5}$. This word may be used to peel the next attacked round.
\subsubsection*{Round 29: input views and the byte loop}
Before this round, recovered subkeys are:
\hx{rk31=9124A012, rk30=01CF72E5}.
The clean view returned by \texttt{peel\_sm4} is
\[C_{\rm view}=\hx{86B3E94F536E42467B938F4C893450AD}.\]
For the clean ciphertext, the peeling calls perform the following actual calculations in order. The last column is the recovered older state word:
\begin{center}\small\begin{tabular}{rllll}\toprule
Undone round & S-box input & $\tau$ output & $L$ output & Older word\\\midrule
31 & 96FF9D45 & 1A488C73 & 138D5078 & 7B938F4C\\
30 & AF8156A0 & 6FBFDA1D & 5B32C6F3 & 893450AD\\
\bottomrule\end{tabular}\end{center}
The first oracle response for this round and its peeled view are
\[\begin{aligned}F&=\hx{FD15CC075271D40B4F8E1DBB536E4246},\\F_{\rm view}&=\hx{4F8E1DBB536E42467B938F4C893450AD}.\end{aligned}\]
The filter operates on these views. In the table, ``fresh'' is the key-byte set for this pair and ``remaining'' is the intersection with earlier observations for the same byte.
\begin{center}\small\begin{tabular}{rrrlll}\toprule
Query & Byte & Mask & $\Delta B$ & Fresh & Remaining\\\midrule
17 & 0 & 01 & 3D000000 & 62 63 & 62 63\\
18 & 0 & 02 & 89000000 & 60 62 & 62\\
19 & 1 & 01 & 00E00000 & 28 29 & 28 29\\
20 & 1 & 02 & 001E0000 & 29 2B & 29\\
21 & 2 & 01 & 0000E900 & 34 35 & 34 35\\
22 & 2 & 02 & 00008200 & 34 36 & 34\\
23 & 3 & 01 & 00000057 & 96 97 & 96 97\\
24 & 3 & 02 & 0000001A & 94 96 & 96\\
\bottomrule\end{tabular}\end{center}
Unique bytes are written into the subkey buffer in the following order. Zero-filled positions not yet written are placeholders, not recovered bytes.
\begin{center}\small\begin{tabular}{rrlr}\toprule
Byte written & Value & Buffer after write & Queries used\\\midrule
0 & 62 & 62000000 & 18\\
1 & 29 & 62290000 & 20\\
2 & 34 & 62293400 & 22\\
3 & 96 & 62293496 & 24\\
\bottomrule\end{tabular}\end{center}
Only now is \texttt{keys[29]} assigned: $rk_{29}=\hx{62293496}$. This word may be used to peel the next attacked round.
\subsubsection*{Round 28: input views and the byte loop}
Before this round, recovered subkeys are:
\hx{rk31=9124A012, rk30=01CF72E5, rk29=62293496}.
The clean view returned by \texttt{peel\_sm4} is
\[C_{\rm view}=\hx{536E42467B938F4C893450AD6FE46B75}.\]
For the clean ciphertext, the peeling calls perform the following actual calculations in order. The last column is the recovered older state word:
\begin{center}\small\begin{tabular}{rllll}\toprule
Undone round & S-box input & $\tau$ output & $L$ output & Older word\\\midrule
31 & 96FF9D45 & 1A488C73 & 138D5078 & 7B938F4C\\
30 & AF8156A0 & 6FBFDA1D & 5B32C6F3 & 893450AD\\
29 & C3E0A931 & 928929B3 & E957823A & 6FE46B75\\
\bottomrule\end{tabular}\end{center}
The first oracle response for this round and its peeled view are
\[\begin{aligned}F&=\hx{84EB762070DAF490AA919B05476A5256},\\F_{\rm view}&=\hx{476A52567B938F4C893450AD6FE46B75}.\end{aligned}\]
The filter operates on these views. In the table, ``fresh'' is the key-byte set for this pair and ``remaining'' is the intersection with earlier observations for the same byte.
\begin{center}\small\begin{tabular}{rrrlll}\toprule
Query & Byte & Mask & $\Delta B$ & Fresh & Remaining\\\midrule
25 & 0 & 01 & 04000000 & 42 43 & 42 43\\
26 & 0 & 02 & 55000000 & 40 42 & 42\\
27 & 1 & 01 & 00820000 & 8C 8D & 8C 8D\\
28 & 1 & 02 & 00450000 & 8D 8F & 8D\\
29 & 2 & 01 & 00005800 & 36 37 & 36 37\\
30 & 2 & 02 & 00006000 & 34 36 & 36\\
31 & 3 & 01 & 00000096 & 54 55 & 54 55\\
32 & 3 & 02 & 00000022 & 54 56 & 54\\
\bottomrule\end{tabular}\end{center}
Unique bytes are written into the subkey buffer in the following order. Zero-filled positions not yet written are placeholders, not recovered bytes.
\begin{center}\small\begin{tabular}{rrlr}\toprule
Byte written & Value & Buffer after write & Queries used\\\midrule
0 & 42 & 42000000 & 26\\
1 & 8D & 428D0000 & 28\\
2 & 36 & 428D3600 & 30\\
3 & 54 & 428D3654 & 32\\
\bottomrule\end{tabular}\end{center}
Only now is \texttt{keys[28]} assigned: $rk_{28}=\hx{428D3654}$. The four required subkeys are now known; leave the round loop and reverse the key schedule.


\subsection{Stage 10: four recovered words enter reverse expansion}
The attack's recovery order was 31, 30, 29, 28. The final list is deliberately
assembled in the ascending order expected by \hx{reverse\_sm4}:
\begin{Verbatim}
final_keys = [0x428D3654, 0x62293496, 0x01CF72E5, 0x9124A012]
\end{Verbatim}
These are $K_{32},K_{33},K_{34},K_{35}$. The first backward calculation is
\[
\begin{aligned}
Q&=\hx{428D3654}\xor\hx{62293496}\xor\hx{01CF72E5}\xor\hx{646B7279}
=\hx{4500025E},\\
\tau(Q)&=\hx{73D6E99D},\qquad T'(Q)=\hx{605CAC93},\\
K_{31}&=\hx{9124A012}\xor\hx{605CAC93}=\hx{F1780C81}.
\end{aligned}
\]
Each following index reuses words already calculated to its right. No new
oracle requests occur in this phase. At the end, remove the four FK words:
\begin{center}\begin{tabular}{rlll}\toprule
$i$ & $K_i$ & $FK_i$ & $MK_i=K_i\oplus FK_i$\\\midrule
0 & A292FFA1 & A3B1BAC6 & 01234567\\
1 & DF01FEBF & 56AA3350 & 89ABCDEF\\
2 & 99A12B0F & 677D9197 & FEDCBA98\\
3 & C42410CC & B27022DC & 76543210\\
\bottomrule\end{tabular}\end{center}
Concatenation returns the 16 bytes \hx{0123456789ABCDEFFEDCBA9876543210}.


\subsection{Stage 11: return values, verification and process exit}
The solver returns \hx{complete=True}, the recovered master key, four subkeys,
32 events, candidate counts and \hx{used=32}. The terminal's schedule display
must agree with that key. The verified round trip is
\[
\begin{aligned}
\operatorname{Encrypt}(P,MK)&=\hx{681EDF34D206965E86B3E94F536E4246}=C,\\
\operatorname{Decrypt}(C,MK)&=\hx{0123456789ABCDEFFEDCBA9876543210}=P.
\end{aligned}
\]
The simulation also compares the recovered key with the target's reference
key, only for validation. \hx{sm4\_demo} returns \hx{True}; \hx{main} returns
0; the logging wrapper finishes and \hx{SystemExit} passes status 0 to the shell.
\note{The 32 faulty ciphertexts are this software example's recovery cost.
Teacher callbacks may recompute values for display; those calculations are not
additional independent observations. Physical glitch calibration and failed
acquisitions are not counted or demonstrated here.}
\clearpage
